%=========== ΛΥΣΗ - 11ο ΘΕΜΑ Δ ===========
\begin{thema}{Δ}
Η $f$ είναι παραγωγίσιμη στο $\mathbb{R}$ και για κάθε $x\in\mathbb{R}$ είναι
\[
f'(x)=3x^2-3=3(x-1)(x+1).
\]
\begin{erwthma}
\item Βρίσκουμε τις ρίζες και το πρόσημο της $f'$:
\begin{itemize}
\item $f'(x)=0\Leftrightarrow x=-1$ ή $x=1$,
\item $f'(x)>0$ για $x<-1$ ή $x>1$,
\item $f'(x)<0$ για $-1<x<1$.
\end{itemize}
Στον παρακάτω πίνακα βλέπουμε το πρόσημο της $f'$ και τη μονοτονία της $f$:
\begin{center}
\begin{tikzpicture}
\tkzTabInit[color,lgt=2,espcl=1.5,colorC = maincolor!50!white,colorL = maincolor!20!white,colorV = maincolor!50!white]%
{$x$ / .8,$f'(x)$ /0.8,$f(x)$/1}%
{$-\infty$,$-1$,$1$,$+\infty$}%
\tkzTabLine{, +, z, -, z, +, }
\tkzTabVar{-/{-\infty},+/{a+2},-/{a-2},+/{+\infty}}
\end{tikzpicture}
\end{center}

Άρα η $f$ είναι \textbf{γνησίως αύξουσα} στα διαστήματα $(-\infty,-1]$ και $[1,+\infty)$ και \textbf{γνησίως φθίνουσα} στο $[-1,1]$.

Επομένως παρουσιάζει:
\begin{itemize}
\item τοπικό μέγιστο στο $x=-1$, με τιμή
\[
f(-1)=-1+3+a=a+2,
\]
\item τοπικό ελάχιστο στο $x=1$, με τιμή
\[
f(1)=1-3+a=a-2.
\]
\end{itemize}

\item Επιπλέον,
\[
\lim_{x\to-\infty}f(x)=-\infty
\quad\text{και}\quad
\lim_{x\to+\infty}f(x)=+\infty.
\]
Για να έχει η εξίσωση $f(x)=0$ ακριβώς τρεις πραγματικές και άνισες μεταξύ τους ρίζες, πρέπει και αρκεί το τοπικό μέγιστο να είναι θετικό και το τοπικό ελάχιστο αρνητικό. Δηλαδή
\begin{align*}
f(-1)>0\quad\text{και}\quad f(1)<0
&\Leftrightarrow a+2>0\quad\text{και}\quad a-2<0\\
&\Leftrightarrow -2<a<2.
\end{align*}
Επομένως
\[
{a\in(-2,2)}.
\]

\item Για $a=2$ είναι
\begin{align*}
f(x)
&=x^3-3x+2\\
&=(x-1)^2(x+2).
\end{align*}
Άρα
\[
f(1)=0
\quad\text{και}\quad
f'(1)=0.
\]
Η εφαπτομένη της $C_f$ στο σημείο $A(1,0)$ έχει εξίσωση
\[
y-f(1)=f'(1)(x-1)\Leftrightarrow {y=0}.
\]
Επομένως η εφαπτομένη στο $A$ είναι ο άξονας $x'x$, άρα η $C_f$ εφάπτεται σε αυτόν.

Για κάθε $x\in\mathbb{R}$ είναι
\[
f''(x)=6x.
\]
Η $f''$ αλλάζει πρόσημο στο $x=0$, οπότε η $C_f$ έχει σημείο καμπής το
\[
K(0,f(0))\equiv K(0,2).
\]
Επειδή
\[
f'(0)=-3,
\]
η εφαπτομένη της $C_f$ στο σημείο καμπής έχει εξίσωση
\begin{align*}
y-f(0)&=f'(0)(x-0)\\
\Leftrightarrow y-2&=-3x\\
\Leftrightarrow {y=-3x+2}.
\end{align*}

\item Για $a=2$ έχουμε
\[
f(x)=(x-1)^2(x+2).
\]
Επομένως
\begin{align*}
\lim_{x\to1}\frac{f(x)-3(x-1)^2}{(x-1)^3}
&=\lim_{x\to1}
\frac{(x-1)^2(x+2)-3(x-1)^2}{(x-1)^3}\\
&=\lim_{x\to1}
\frac{(x-1)^2(x-1)}{(x-1)^3}\\
&={1}.
\end{align*}
\end{erwthma}
\end{thema}
%=========== ΤΕΛΟΣ ΛΥΣΗΣ - 11ο ΘΕΜΑ Δ ===========
